\documentclass[12pt,a4paper]{article}

\usepackage[round,authoryear]{natbib}
\usepackage{amsmath}
\usepackage{physics}

\title{Introduction to \LaTeX}
\author{Shlok Patel}
\date{\today}

\begin{document}

\maketitle

% Task 1: Hello World
\section{Hello World}
Hello World!

% Task 2: first document and table of contents
\section{My First \LaTeX\ Document}
This is a paragraph. Multiple spaces in the source are generally
treated as one space in the output. A blank line starts a new paragraph.

Special characters must be escaped:
\# \quad \$ \quad \% \quad \& \quad \_ \quad
\{ \quad \} \quad \textbackslash.

\subsection{A Subsection}
This is a subsection. Add more sections with \verb|\section|
and \verb|\subsection|.

\subsection{Numbered List}
\begin{enumerate}
    \item First item
    \item Second item
    \item Third item
\end{enumerate}

% The table of contents is placed near the beginning of the document.
\clearpage
\tableofcontents
\clearpage

% Task 2.9
\section{Science!}
Einstein's mass-energy equivalence relation is
\(E=mc^2\).

% Task 3: citations, equation and cross-reference
\section{Equations, Citations and Cross-References}

Einstein's mass-energy equivalence paper is cited here
\citep{einstein1905}. The famous relation is \(E=mc^2\).

The time-dependent Schrödinger equation can be written as
\begin{equation}
i\hbar\frac{\partial}{\partial t}\Psi(\mathbf{r},t)
=
\left[
-\frac{\hbar^2}{2\mu}\nabla^2+V(\mathbf{r},t)
\right]\Psi(\mathbf{r},t).
\label{eq:schrodinger}
\end{equation}

Equation~\eqref{eq:schrodinger} is the Schrödinger equation.
The original 1926 eigenvalue-problem paper is cited as
\citet{schrodinger1926}.

\bibliographystyle{agu}
\bibliography{physics}

\end{document}
